\begin{axiom} \label{an1:ax:peano}
  The set of \textbf{natural numbers} $\N$ satisfies the following axioms:
  \begin{itemize}
    \item $1 \in \N$
    \item $\forall n \in \N, s(n) \in \N$
    \item $\nexists n \in \N, s(n) = 1$
    \item $\forall m, n \in \N, s(m) = s(n) \implies m = n$
    \item $\forall S \subseteq \N, (1 \in S \land (\forall n \in S, s(n) \in S)) \implies S = \N$
  \end{itemize}
  $s(n)$ is called the \textbf{successor} of $n$. These axioms are called the \textbf{Peano axioms}.
\end{axiom}

\begin{remark}
  The \textbf{Peano axioms} do not guarantee the \emph{existence} of $\N$, which is guaranteed by the
  \textbf{axiom of infinity} in the \textbf{ZFC axioms}. In this textbook, the \textbf{ZFC axioms}
  are not discussed, so we will just assume that $\N$ exists.
\end{remark}

\begin{theorem} \label{an1:thm:mathematical-induction}
  Suppose $P$ is a predicate with domain $\N$. Then
  \[
    P(1) \land (\forall n \in \N, P(n) \implies P(s(n))) \implies \forall n \in \N, P(n)
  \]
  This method of proof is called the \textbf{principle of mathematical induction}.
\end{theorem}

\begin{proof}
  Suppose
  \begin{itemize}
    \item $P(1)$
    \item $\forall n \in \N, P(n) \implies P(s(n))$
  \end{itemize}
  Define
  \[
    S \coloneqq \setbuilder*{n \in \N}{P(n)} \subseteq \N
  \]
  Then
  \begin{itemize}
    \item $1 \in S$
    \item $\forall n \in S, s(n) \in S$
  \end{itemize}
  By \cref{an1:ax:peano},
  \[
    \begin{array}{c}
      S = \N     \\
      \Downarrow \\
      \forall n \in \N, P(n)
    \end{array}
  \]
\end{proof}

\begin{definition} \label{an1:def:operations-on-N}
  Suppose $n, m \in \N$.
  \begin{itemize}
    \item \textbf{Addition} $+$ on $\N$ is defined by
          \begin{itemize}
            \item $n + 1 \coloneqq s(n)$
            \item $n + s(m) \coloneqq s(n + m)$
          \end{itemize}
    \item \textbf{Multiplication} $\cdot$ on $\N$ is defined by
          \begin{itemize}
            \item $n \cdot 1 \coloneqq n$
            \item $n \cdot s(m) \coloneqq n \cdot m + n$
          \end{itemize}
  \end{itemize}
\end{definition}

\begin{remark}
  Strictly speaking, \cref{an1:def:operations-on-N} requires the \textbf{recursion theorem} to be strictly defined.
  In this textbook, the \textbf{recursion theorem} is not discussed, so we will just assume that the operations
  on $\N$ are \textbf{well-defined}.
\end{remark}

\begin{theorem} \label{an1:thm:well-ordering-property}
  $\N$ has the \textbf{well-ordering property}, which states
  \[
    \forall S \subseteq \N, S \neq \vnot \implies \exists m \in S, m = \min S
  \]
\end{theorem}

\begin{proof}
  Assume
  \[
    \exists S \subseteq \N, S \neq \vnot \land (\nexists m \in S, m = \min S)
  \]
  Define
  \begin{itemize}
    \item $T \coloneqq \N \setminus S \subseteq \N$
    \item $L \coloneqq \setbuilder*{n \in \N}{\set*{1, 2, \dots, n} \subseteq T} \subseteq \N$
  \end{itemize}
  We will show
  \[
    \begin{array}{c}
      L = \N     \\
      \Downarrow \\
      T = \N     \\
      \Downarrow \\
      S = \vnot  \\
      \Downarrow \\
      \bot
    \end{array}
  \]
  \begin{itemize}
    \item \textbf{Base case:} Assume $1 \in S$. Then
          \[
            \begin{array}{c}
              1 = \min S \\
              \Downarrow \\
              \bot
            \end{array}
          \]
          Therefore,
          \[
            \begin{array}{c}
              1 \notin S \\
              \Downarrow \\
              1 \in T    \\
              \Downarrow \\
              1 \in L
            \end{array}
          \]
    \item \textbf{Inductive step:} Suppose $n \in L$. Then
          \[
            \begin{array}{c}
              n \in L                           \\
              \Downarrow                        \\
              \set*{1, 2, \dots, n} \subseteq T \\
              \Downarrow                        \\
              \set*{1, 2, \dots, n} \cap S = \vnot
            \end{array}
          \]
          Assume $n + 1 \in S$. Then
          \[
            \begin{array}{c}
              n + 1 = \min S \\
              \Downarrow     \\
              \bot
            \end{array}
          \]
          Therefore,
          \[
            \begin{array}{c}
              n + 1 \notin S                        \\
              \Downarrow                            \\
              n + 1 \in T                           \\
              \Downarrow                            \\
              \set*{1, 2, \dots, n + 1} \subseteq T \\
              \Downarrow                            \\
              n + 1 \in L
            \end{array}
          \]
  \end{itemize}
  By \cref{an1:thm:mathematical-induction},
  \[
    L = \N
  \]
\end{proof}
